\section{Autonomy}
\label{sec:autonomy}

\subsection{Mission and Architecture}
The mission has four parts.
\begin{enumerate}
  \item Follow the line from the start to the start square.
  \item Visit the three drop zones and read at each the colour that it accepts.
  \item For each pick bay in turn: approach it, read the colour of its coin, cage the coin,
        carry it to the drop zone of that colour and release it.
  \item Throughout, drive around obstacle boxes that are not in the map.
\end{enumerate}

The software is a set of ROS~2 packages (\cref{fig:architecture}). Odometry, mapping,
localisation, path planning and path tracking are taken from existing open-source packages
and configured for this robot. The line follower, the docking sequences, the mission state
machine and the referee are written for it (\cref{tab:sources}). The packages for
navigation and for the mission have no knowledge that they run in a simulator: everything
that exists only in simulation, including the emulation of the infrared array and the
referee, is kept in a separate package.

\begin{table}[htbp]
  \centering
  \caption{Origin of the components of the autonomy stack.}
  \label{tab:sources}
  \small
  \begin{tabular}{@{}lll@{}}
    \toprule
    Function & Solution & Source \\
    \midrule
    Odometry & filter of wheel speed and yaw rate & \code{robot\_localization}~\cite{robotlocalization} \\
    Mapping, localisation & RGB-D graph SLAM & RTAB-Map~\cite{rtabmap} \\
    Planning and tracking & NavFn, Regulated Pure Pursuit & Nav2~\cite{nav2} \\
    QR decoding & bar code reader on the colour image & ZBar~\cite{zbar} \\
    Line following & PD steering on the line offset & this work \\
    Pick and drop sequences & infrared array and odometry & this work \\
    Mission sequencing & state machine & this work \\
    Scoring & referee on simulator ground truth & this work \\
    \bottomrule
  \end{tabular}
\end{table}

\begin{figure}[htbp]
  \centering
  \begin{tikzpicture}
    \node[ext, minimum width=24mm, minimum height=74mm] (sim) at (0,0)
      {Gazebo\\[3pt]\scriptsize robot, arena,\\\scriptsize sensors};
    \node[box, minimum width=30mm] (ekf) at (6.1,3.0) {odometry filter};
    \node[box, minimum width=30mm] (slam) at (6.1,1.5) {RTAB-Map};
    \node[box, minimum width=30mm] (nav) at (6.1,0.0) {Nav2};
    \node[box, minimum width=30mm] (line) at (6.1,-1.5) {line follower};
    \node[box, minimum width=30mm] (qr) at (6.1,-3.0) {QR reader};
    \node[box, minimum width=28mm, minimum height=50mm] (mis) at (11.6,-1.2)
      {mission\\[3pt]\scriptsize state machine,\\\scriptsize docking};
    \node[ext, minimum width=28mm] (ref) at (11.6,3.0) {referee};
    \draw[->] (sim.east |- ekf) -- node[note, above] {wheel odometry, IMU} (ekf);
    \draw[->] (sim.east |- slam) -- node[note, above] {colour, depth} (slam);
    \draw[->] (sim.east |- nav) -- node[note, above] {depth cloud} (nav);
    \draw[->] (sim.east |- line) -- node[note, above] {IR voltages} (line);
    \draw[->] (sim.east |- qr) -- node[note, above] {colour image} (qr);
    \draw[->] (ekf) -- node[note, right] {odometry} (slam);
    \draw[->] (slam) -- node[note, right] {map, pose} (nav);
    \draw[->] (mis.west |- nav) -- node[note, above] {goal} (nav);
    \draw[->] (mis.west |- line) -- node[note, above] {goal} (line);
    \draw[->] (qr) -- node[note, above] {colour} (mis.west |- qr);
    \draw[->, dashed] (mis) -- node[note, right] {state} (ref);
    \draw[->] (mis.south) -- ++(0,-0.75) -| node[note, pos=0.25, below] {velocity command, jaw angle} (sim.south);
    \draw[->, dashed] (sim.north) -- ++(0,0.55) -| node[note, pos=0.25, above] {true poses of the coins} (ref.north);
  \end{tikzpicture}
  \caption{Software architecture. Solid boxes run unchanged on a real robot; dashed boxes
  exist only in simulation. Velocity commands are issued by Nav2, by the line follower or by
  the mission itself, one at a time, because the mission waits for each step to finish before
  it starts the next.}
  \label{fig:architecture}
\end{figure}

\subsection{Odometry}
\label{sec:odometry}
An extended Kalman filter in planar mode fuses the forward speed from the wheel encoders
with the yaw rate from the gyroscope, at 30\unit{Hz}. The yaw rate of the wheel odometry is
deliberately left out: when the robot turns on the spot the wheels slip, and the gyroscope
is the better source. The filter publishes the transform from the odometry frame to the
robot.

\subsection{Mapping and Localisation}
\label{sec:slam}
\paragraph{Building the map.} The map is built once, in the empty arena. A survey node
sends a fixed route of seven stops to the navigation and turns the robot through a full
circle at each stop, so that the camera sees every wall, while RTAB-Map adds a node to its
graph every 10\unit{cm} or 0.2\unit{rad} of travel. The route ends at the start pose and
closes the loop there. Because the survey starts at the start pose, the map frame coincides
with the frame of the arena, and the station poses can be written in arena coordinates. The
map is constrained to the plane, and its occupancy grid has cells of 2\unit{cm}. Points
below 3\unit{cm} are treated as floor, which keeps coins out of the map.

\paragraph{Localising in it.} A mission run loads the whole map and starts at its origin,
where the robot stands. Each time RTAB-Map recognises a view, twice a second at most, it
corrects the transform from the map to the odometry frame. Between corrections the pose is
carried by the odometry filter.

\paragraph{Look-alike views.} The arena contains the same QR code in two places, once at a
pick bay and once at a drop zone. Seen from close by, a code fills much of the image, and a
match of the features of the code alone was enough for RTAB-Map to conclude that the robot
stood in front of the other wall. The pose then jumped across the room. The remedy is a
condition on how the matched features are spread over the image: a view of the room matches
across the whole width of the walls, a view of a code only within the code. The walls fill a
low band of the image, and during development real views scored between 0.004 and 0.014 on
this measure and look-alikes below 0.0025; the threshold is set at 0.005. The choice is a
compromise. A threshold at the lower end of the real range rejects some true matches, and
one set higher rejects them all, after which the robot runs on odometry alone. In addition,
every new fix is checked against the last fifty views and the odometry between them.

\subsection{Navigation}
\label{sec:navigation}
Between stations the robot is driven by Nav2. The global planner is NavFn on a costmap of
2\unit{cm} cells built from the saved map, and the path is planned again once a second,
which is what steers the robot around boxes that are not in the map. The path is tracked by
the Regulated Pure Pursuit controller at 0.15\unit{m/s} with a look-ahead of 0.20\unit{m},
and the robot is turned on the spot when the path departs from its heading by more than
0.5\unit{rad}. A goal counts as reached within 2\unit{cm} and 0.05\unit{rad}.

Obstacles enter the costmaps from the depth image, reduced to every fourth pixel and to
cells of 2\unit{cm}. Points between 3\unit{cm} and 0.40\unit{m} above the floor mark
obstacles, and rays through free space clear them. The robot is treated as a disc of
0.13\unit{m} radius, and obstacles are inflated by 0.25\unit{m}.

Three properties of the robot shape the configuration.
\begin{itemize}
  \item The camera looks forward only. The controller is not allowed to reverse, and the
        recovery that backs the robot up is removed from the behaviour tree. The remaining
        recoveries clear the costmaps, turn the robot through a quarter circle to look
        around, and wait.
  \item Low objects do not reach the costmap. This is intended for coins and sensor noise,
        and it means that an obstacle lower than 3\unit{cm} would not be avoided.
  \item Turns on the spot are slow, 0.3 to 0.6\unit{rad/s}, so that successive camera views
        overlap and the localisation can follow.
\end{itemize}

\subsection{Line Follower}
\label{sec:linefollower}
\paragraph{Line position.} Each voltage $V_i$ is converted to a darkness
$d_i=1-V_i/V_{\mathrm{ref}}$ between 0 on white and 1 on black. Readings below 0.15 are set
to zero, so that noise on the white floor carries no weight. The offset of the line from the
centre of the array is the weighted mean
\begin{equation}
  e = p\,\frac{\sum_i d_i\,(i-3.5)}{\sum_i d_i}, \qquad i=0,\dots,7,
  \label{eq:offset}
\end{equation}
with $p$ the pitch. It is positive when the line is to the left. Because a line narrower
than the pitch darkens two neighbouring channels partially (\cref{sec:markings}), $e$ varies
continuously as the line moves across the array, and the resolution is much finer than the
pitch.

\paragraph{Steering.} The robot drives at a constant speed $v$ and steers with
\begin{equation}
  \omega = \operatorname{sat}_{\omega_{\max}}\bigl(k_p\,e + k_d\,\dot e\bigr),
  \label{eq:steer}
\end{equation}
with $k_p=45\unit{rad/s}$ per metre, $k_d=2\unit{rad/s}$ per metre per second and
$\omega_{\max}=1.5\unit{rad/s}$, at 20\unit{Hz}. If the line leaves the array, the robot
turns at 0.6\unit{rad/s} towards the side on which the line was last seen. If no line has
been seen yet, it drives straight, so that a goal may start a little short of a stem.

\paragraph{Why the position of the array matters.} Let the array be a distance $L$ ahead of
the axle, let $\psi$ be the heading of the robot relative to a straight line, and let $y$ be
the lateral position of the array relative to the line, so that $e=-y$. For small angles,
\begin{equation}
  \dot y = v\,\psi + L\,\omega, \qquad \dot\psi = \omega .
\end{equation}
With the proportional part of \cref{eq:steer}, $\omega=-k_p\,y$, this gives
\begin{equation}
  \ddot y + L\,k_p\,\dot y + v\,k_p\,y = 0,
  \label{eq:line-dynamics}
\end{equation}
a second-order system with natural frequency $\omega_n=\sqrt{v\,k_p}$ and damping ratio
$\zeta=\tfrac{1}{2}L\sqrt{k_p/v}$. The damping comes entirely from $L$: an array on the axle
($L=0$) would give an undamped oscillation under proportional control, whatever the gain.
With $L=48\unit{mm}$ and $k_p=45$, the damping ratio is 0.51 at the entry speed of
0.10\unit{m/s} and 0.72 at the docking speed of 0.05\unit{m/s}. The derivative term raises
the first value slightly, to 0.53. The mechanical decision to hang the array ahead of the
axle is thus what makes a simple controller sufficient, and driving the stems slowly is what
makes the approach to a marker well damped.

\paragraph{Markers.} A marker is reported when at least six channels have a darkness of 0.5
or more in two successive samples. A line darkens two channels at most, so only a square or
a crossbar can do this, and the requirement of two samples rejects a single noisy reading.
Two sample periods at 20\unit{Hz} are 5\unit{mm} of travel at the docking speed.

The follower is offered as an action. A goal names a maximum distance and a speed, and the
result states whether a marker was reached and how far the robot drove.

\subsection{Docking}
\label{sec:docking}
A station is approached in two stages (\cref{fig:docking}). Nav2 brings the robot to the
station pose. The line follower then drives straight until the array finds the stem, steers
onto it, and stops at the marker. Everything after that is a timed movement of the jaws or
a straight move measured by odometry.

\begin{description}
  \item[Pick.] Open the jaws. Follow the stem to the crossbar, at which point the coin is
        between the jaws. Read the QR code on the wall. Close the jaws. Back out by the
        distance just driven.
  \item[Drop.] Follow the stem until the array enters the square. Open the jaws, which
        pushes the coin clear. Back out, then close the jaws.
  \item[Read.] Follow the stem to its marker, read the QR code and back out. This is used
        for the survey of the drop zones.
\end{description}
The robot backs out along the stem that it has just driven and by the distance that it
measured, which is the only reverse movement in the mission. It needs no sensing behind the
robot because the path has just been traversed.

\begin{figure}[htbp]
  \centering
  \begin{tikzpicture}[x=9cm, y=9cm]
    % pick bay drawn with the approach direction to the right (x along the stem)
    \draw[draw=ink, line width=1.1pt] (0.95,-0.17) -- (0.95,0.17);
    \node[note, rotate=-90, anchor=south] at (0.96,0) {wall with QR code};
    \draw[draw=black, line width=1.2pt] (0.30,0) -- (0.60,0);
    \draw[draw=black, line width=1.2pt] (0.60,-0.06) -- (0.60,0.06);
    \fill[cB] (0.73,0) circle (0.02);
    \node[note, anchor=south, text=cB] at (0.73,0.03) {coin};
    % station pose with a lateral error
    \draw[draw=ink, fill=panel] (0.133,-0.045) rectangle (0.313,0.105);
    \fill[ink] (0.15,0.03) circle (0.006);
    \draw[->, cA, line width=0.8pt] (0.15,0.03) .. controls (0.30,0.03) and (0.33,0.0) .. (0.42,0.0) -- (0.545,0.0);
    \node[note, anchor=north] at (0.22,-0.05) {robot at the station pose};
    \node[note, anchor=south, text=cA] at (0.44,0.06) {line follower};
    \draw[|-|] (0.15,-0.13) -- (0.30,-0.13);
    \node[note, below] at (0.225,-0.135) {0.15\,m};
    \draw[|-|] (0.30,-0.13) -- (0.60,-0.13);
    \node[note, below] at (0.45,-0.135) {stem, 0.30\,m};
    \draw[|-|] (0.60,-0.13) -- (0.73,-0.13);
    \node[note, below] at (0.665,-0.135) {0.13\,m};
    \node[note, anchor=south] at (0.60,0.07) {crossbar};
  \end{tikzpicture}
  \caption{Approach to a pick bay, to scale. The navigation leaves the robot at the station
  pose with an error of a few centimetres; the line follower removes it along the stem and
  stops the array on the crossbar, with the coin between the open jaws.}
  \label{fig:docking}
\end{figure}

\paragraph{QR codes are read at the station only.} The decoder reports the text of a code
but not where the code is in the image. A code is therefore trusted only when the robot
stands at the marker in front of it, where no neighbouring code is in view, and only after
three identical decodes in a row. Decodes received before the robot arrived are discarded,
so that a code passed on the way is never taken for the code of the station.

\subsection{Mission State Machine}
\label{sec:statemachine}
The mission is a sequence of blocking steps, and the state machine only decides what happens
next (\cref{fig:states}). Its failure handling is deliberately plain.
\begin{itemize}
  \item If the start square is not found within 3\unit{m} of travel, the mission fails.
  \item If a marker is not found at a station, the robot backs out and tries the station once
        more.
  \item A drop zone whose code cannot be read is left out of the table of colours.
  \item A coin whose code cannot be read is left in its bay, with the jaws still open, since
        without a colour there is nowhere to take it.
  \item A coin whose colour has no known drop zone is set down where the robot stands.
\end{itemize}
The mission reports success only when all three coins have been delivered.

\begin{figure}[htbp]
  \centering
  \begin{tikzpicture}
    \node[box, minimum width=27mm] (a) at (0,0) {follow\\entry line};
    \node[box, minimum width=27mm] (b) at (3.6,0) {survey\\drop zones};
    \node[box, minimum width=27mm] (c) at (7.2,0) {pick\\bay $k$};
    \node[box, minimum width=27mm] (d) at (10.8,0) {drop\\by colour};
    \node[box, minimum width=18mm] (e) at (11.3,-1.9) {done};
    \node[box, minimum width=20mm] (f) at (0,-1.9) {failed};
    \draw[->] (a) -- (b); \draw[->] (b) -- (c); \draw[->] (c) -- (d);
    \draw[->] (d.north) -- ++(0,0.5) -| node[note, pos=0.25, above] {next bay} (c.north);
    \draw[->] ([xshift=5mm]d.south) -- node[note, right] {three delivered} (e.north);
    \draw[->] (a) -- node[note, right] {no start square} (f);
    \draw[->] ([xshift=-12mm]d.south) |- node[note, pos=0.75, above] {fewer than three delivered} (f.east);
  \end{tikzpicture}
  \caption{States of the mission.}
  \label{fig:states}
\end{figure}

\subsection{Referee}
\label{sec:referee}
A run is scored by a referee that reads the true positions of the coins from the simulator.
A coin is home when its centre lies inside the drop square whose QR code names its colour.
The referee has no connection to the robot other than the state of the mission, which tells
it when to give its verdict, and nothing on the robot may read the ground truth. A run
passes when all three coins are home at the end.
